% Answers to Chapter 5, from lectures/L05/appendix/c_solutions.md; the self-test from
% lectures/L05/README.md.

\facitavsnitt{c5}{Kapitel 5}{Potenser och rötter}
\begin{facit}

\svar{1.1} $I_1 = 2\,\text{A}$, $I_2 = 3\,\text{A}$

\svar{1.2} $U_A = \frac{54}{11} \approx 4{,}91\,\text{V}$, $U_B = \frac{30}{11} \approx
2{,}73\,\text{V}$

\svar{2.1} \sv{a} $x^8$ \sv{b} $R^4$ \sv{c} $I^6$ \sv{d} $U^3$ \sv{e} $\frac{8}{R^3}$

\svar{2.2} \sv{a} $4{,}7 \times 10^4\,\Omega$ \sv{b} $2{,}2 \times 10^{-5}\,\text{F}$
\sv{c} $3{,}3 \times 10^6\,\text{Hz}$ \sv{d} $1{,}5 \times 10^{-8}\,\text{A}$

\svar{2.3} \sv{a} $250\,\text{mW}$ \sv{b} $80\,\text{mW}$

\svar{2.4} \sv{a} $\approx 230\,\text{V}$ \sv{b} $\approx 17\,\text{V}$

\svar{3.1} \sv{a} $81$ \sv{b} $-32$ \sv{c} $0{,}01$ \sv{d} $1$ \sv{e} $\frac{1}{8} = 0{,}125$
\sv{f} $0{,}001$

\svar{3.2} \sv{a} $a^4$ \sv{b} $\frac{1}{U^3}$ \sv{c} $4R^6$ \sv{d} $I^5$ \sv{e} $\frac{R^6}{27}$
\sv{f} $3x^3$

\svar{3.3} \sv{a} $10^{-3}$ \sv{b} $10^3$ \sv{c} $10^{-6}$ \sv{d} $10^7$ \sv{e} $10^{-3}$

\svar{3.4} \sv{a} $9$ \sv{b} $0{,}5$ \sv{c} $3$ \sv{d} $0{,}04$ \sv{e} $5\sqrt{2}$
\sv{f} $\sqrt{9 + 16} = \sqrt{25} = 5$, men $\sqrt{9} + \sqrt{16} = 3 + 4 = 7$.

\svar{3.5} \sv{a} $4$ \sv{b} $3$ \sv{c} $27$ \sv{d} $\frac{1}{2}$ \sv{e} $2$

\svar{3.6} \sv{a} $2{,}2 \times 10^3\,\Omega$ \sv{b} $4{,}7 \times 10^{-4}\,\text{F}$
\sv{c} $1{,}5 \times 10^{-11}\,\text{F}$ \sv{d} $2{,}5 \times 10^{-1}\,\text{A}$
\sv{e} $1{,}8 \times 10^6\,\text{Hz}$ \sv{f} $6{,}8 \times 10^{-9}\,\text{H}$

\svar{3.7} \sv{a} $3{,}3\,\mu\text{F}$ \sv{b} $470\,\text{k}\Omega$ \sv{c} $250\,\mu\text{A}$
\sv{d} $80\,\text{MHz}$ \sv{e} $12\,\text{pF}$

\svar{3.8} \sv{a} $8 \times 10^{-3}$ \sv{b} $3 \times 10^{-5}$ \sv{c} $2{,}5 \times 10^7$
\sv{d} $3{,}25 \times 10^{-3}$ \sv{e} $9 \times 10^1$

\svar{3.9} \sv{a} Tre. \sv{b} $0{,}025$ \sv{c} $1{,}23 \times 10^4\,\Omega$
\sv{d} Två värdesiffror, som spänningen: $I = 0{,}050\,\text{A} = 50\,\text{mA}$.

\svar{3.10} \sv{a} $15\,\text{mA}$ \sv{b} $11{,}75\,\text{V}$ \sv{c} $3\,\text{k}\Omega$

\svar{3.11} \sv{a} $0{,}6\,\text{W}$ \sv{b} $\approx 11{,}4\,\text{mW}$ \sv{c} $\approx
23\,\text{mA}$ \sv{d} $\approx 10{,}8\,\text{V}$

\svar{3.12} \sv{a} $47\,\text{ms}$ \sv{b} $10\,\text{k}\Omega$ \sv{c} $500\,\text{nF}$

\svar{3.13} \sv{a} $\approx 7{,}07\,\text{V}$ \sv{b} $\approx 325\,\text{V}$
\sv{c} $U_{\text{RMS}}^2 = \frac{\abs{U}^2}{(\sqrt{2})^2} = \frac{\abs{U}^2}{2}$, så
$P = \frac{\abs{U}^2}{2R}$.
\sv{d} $9\,\text{W}$

\svar{3.14} \sv{a} Exponenterna ska adderas: $2^7 = 128$.
\sv{b} Exponenterna ska multipliceras: $a^6$.
\sv{c} Exponenterna ska subtraheras: $10^4$.
\sv{d} Exponenten gäller båda faktorerna: $8a^3$.
\sv{e} En rot kan inte delas upp över en summa: $\sqrt{25} = 5$.
\sv{f} En negativ exponent ger en reciprok, inte ett negativt tal: $\frac{1}{a^2}$.

\svar[Testa dig själv]{T} \sv{1} $3x^2$ \sv{2} $2{,}2 \times 10^{-7}\,\text{F}$

\end{facit}
